% !TEX program = xelatex
% callouts: ready-made kinds, \akizumiNewCallout, \optional, numbering, cref,
% page breaks, and a preamble's own definitions winning
\documentclass{akizumi-book}
\title{Callouts}
\akizumiNewCallout{keypoint}{Key point}{學習要點}{}{}[numbered=chapter, tint]
\akizumiNewCallout{stylepoint}{Style point}{風格要點}{}{}[numbered, color=asidebar, labelcolor=ink]
\akizumiCalloutSetup{keyidea}{numbered=chapter}
\newcommand{\pitfall}[1]{\textbf{Pitfall.}~#1}  % user's own: the env is skipped
\begin{document}
\chapter{Series}
\begin{chapterguide}
\guidehead{What this chapter is about}
When does an infinite sum have a value?
\guidehead{Prerequisites}
Limits of sequences.
\end{chapterguide}

\begin{keyidea}\label{ki:tails}
A series converges when its tails become small.
\end{keyidea}

\begin{theorem}[Comparison]\label{thm:cmp}
If $\abs{a_n}\le b_n$ and $\sum b_n$ converges, so does $\sum a_n$.
\end{theorem}
\begin{proofidea}
Tails of $\sum\abs{a_n}$ are bounded by tails of $\sum b_n$.
\end{proofidea}
\begin{proof}
Cauchy criterion.
\end{proof}

\begin{tip}[Root test]
Try it when $a_n$ has an $n$-th power.
\end{tip}
\begin{warning}
Rearranging a conditionally convergent series can change its sum.
\end{warning}
\pitfall{Here \texttt{\textbackslash pitfall} is the preamble's command.}
\begin{note}
Notes are quiet asides.
\end{note}
\begin{setup}
Let $(a_n)$ be a sequence of real numbers and $s_n=\sum_{k\le n}a_k$.
\end{setup}
\begin{setup}[Data]
A titled setup box shows only its title.
\end{setup}
\begin{keypoint}\label{kp:one}
你不需要在 \code{import} 之前先確認套件裝了沒。
\end{keypoint}
\begin{stylepoint}
Name things after what they mean.
\end{stylepoint}
See \cref{ki:tails}, \cref{kp:one} and \cref{thm:cmp}.

\section{\optional Rearrangements}\label{sec:opt}
\begin{remark}[\optional Riemann]
Optional remark.
\end{remark}
\begin{keypoint}
Second key point: \thekeypoint.
\end{keypoint}

\begin{keyidea}[Breaks across pages]
\foreach \i in {1,...,60}{Line \i\ of a long key idea that must split across pages.\par}
\end{keyidea}

\chapter{Next}
\begin{keypoint}\label{kp:two}
Numbering restarts: \thekeypoint.
\end{keypoint}
\end{document}
